% Options for packages loaded elsewhere
\PassOptionsToPackage{unicode}{hyperref}
\PassOptionsToPackage{hyphens}{url}
\documentclass[
  11pt,
]{article}
\usepackage{xcolor}
\usepackage[margin=2.5cm]{geometry}
\usepackage{amsmath,amssymb}
\setcounter{secnumdepth}{-\maxdimen} % remove section numbering
\usepackage{iftex}
\ifPDFTeX
  \usepackage[T1]{fontenc}
  \usepackage[utf8]{inputenc}
  \usepackage{textcomp} % provide euro and other symbols
\else % if luatex or xetex
  \usepackage{unicode-math} % this also loads fontspec
  \defaultfontfeatures{Scale=MatchLowercase}
  \defaultfontfeatures[\rmfamily]{Ligatures=TeX,Scale=1}
\fi
\usepackage{lmodern}
\ifPDFTeX\else
  % xetex/luatex font selection
\fi
% Use upquote if available, for straight quotes in verbatim environments
\IfFileExists{upquote.sty}{\usepackage{upquote}}{}
\IfFileExists{microtype.sty}{% use microtype if available
  \usepackage[]{microtype}
  \UseMicrotypeSet[protrusion]{basicmath} % disable protrusion for tt fonts
}{}
\makeatletter
\@ifundefined{KOMAClassName}{% if non-KOMA class
  \IfFileExists{parskip.sty}{%
    \usepackage{parskip}
  }{% else
    \setlength{\parindent}{0pt}
    \setlength{\parskip}{6pt plus 2pt minus 1pt}}
}{% if KOMA class
  \KOMAoptions{parskip=half}}
\makeatother
\usepackage{longtable,booktabs,array}
\usepackage{calc} % for calculating minipage widths
% Correct order of tables after \paragraph or \subparagraph
\usepackage{etoolbox}
\makeatletter
\patchcmd\longtable{\par}{\if@noskipsec\mbox{}\fi\par}{}{}
\makeatother
% Allow footnotes in longtable head/foot
\IfFileExists{footnotehyper.sty}{\usepackage{footnotehyper}}{\usepackage{footnote}}
\makesavenoteenv{longtable}
\setlength{\emergencystretch}{3em} % prevent overfull lines
\providecommand{\tightlist}{%
  \setlength{\itemsep}{0pt}\setlength{\parskip}{0pt}}
\usepackage{bookmark}
\IfFileExists{xurl.sty}{\usepackage{xurl}}{} % add URL line breaks if available
\urlstyle{same}
\hypersetup{
  pdftitle={Dosimetric-Response Modeling in Glioblastoma (CFB-GBM)},
  pdfauthor={TCP Modeling GBM Project},
  hidelinks,
  pdfcreator={LaTeX via pandoc}}

\title{Dosimetric-Response Modeling in Glioblastoma (CFB-GBM)}
\author{TCP Modeling GBM Project}
\date{}

\begin{document}
\maketitle

\section{Dosimetric--Response Modeling in Glioblastoma Using the Open
CFB-GBM Cohort: A TCP Pipeline Feasibility
Study}\label{dosimetricresponse-modeling-in-glioblastoma-using-the-open-cfb-gbm-cohort-a-tcp-pipeline-feasibility-study}

\textbf{Draft manuscript} --- verified results from
\texttt{reports/RESULTS.md} (regenerate: \texttt{make\ report}).\\
Reference map: \texttt{reports/literature\_table.csv} (ref\_id 1--18).

\begin{center}\rule{0.5\linewidth}{0.5pt}\end{center}

\subsection{Abstract}\label{abstract}

\textbf{Background:} Tumor control probability (TCP) models relate
radiotherapy dose to local tumour control {[}4,5{]}. Open imaging
cohorts promise reproducible validation, but routine clinical data often
lack classical TCP endpoints and dose heterogeneity {[}7{]}.

\textbf{Methods:} We built a reproducible Python pipeline on the TCIA
CFB-GBM cohort (n=190 after DVH quality control) {[}8{]}. Four TCP
models (Poisson, logistic, probit, gEUD) {[}4--6{]} were fit with
bootstrap confidence intervals {[}18{]}. Version 3 supplementary data
(RANO imaging response {[}12,13{]}, n=137 with t0→t1 labels) were
integrated {[}8,9{]}. We audited confounding, compared OS-based vs
RANO-based endpoints, and performed within-arm analyses stratified by
fractionation scheme (60 Gy/30 fr vs 40.05 Gy/15 fr) {[}2,3,15{]}.

\textbf{Results:} Pooled EQD2 discriminated an OS median-split proxy
(ROC AUC≈0.68) but not RANO non-progressive disease when used as a
dose-only TCP input (AUC≈0.43). After adjusting for fractionation
scheme, \textbf{pooled GTV volume + clinical covariates achieved
AUC≈0.72} (n=137); LOOCV AUC≈0.64. Within the hypofractionated arm
(n=34) {[}3,15{]}, pre-treatment GTV volume predicted RANO non-PD
(multivariable in-sample AUC=0.90; \textbf{LOOCV AUC=0.74}). PyRadiomics
features (t1gd GTV t0, v3 TSV) {[}9,17{]} outperformed DVH volume alone
for RANO (top-5 in-sample AUC≈0.78 vs 0.71; \textbf{nested 5-fold CV
AUC≈0.74 vs 0.70}). GTV volume from DVH matched RANO t0 segmentation
volumes (Spearman ρ=1.00, n=141) {[}8{]}.

\textbf{Conclusions:} Classical pooled TCP dose--response validation is
not supported on CFB-GBM; however, \textbf{tumour burden models} (volume
and radiomics) predict early RANO when fractionation confounding is
addressed. We provide an open pipeline, LOOCV benchmarks, and a
feasibility checklist for TCP studies on routine RT datasets {[}7{]}.

\textbf{Keywords:} glioblastoma; tumor control probability; DVH; RANO;
open data; CFB-GBM

\begin{center}\rule{0.5\linewidth}{0.5pt}\end{center}

\subsection{1. Introduction}\label{introduction}

Glioblastoma remains the most common malignant primary brain tumour in
adults. Standard treatment includes maximal safe resection followed by
radiotherapy with concurrent temozolomide {[}1{]}. Despite uniform
protocols, outcome varies widely, motivating dose--response and imaging
biomarker research {[}2,9{]}.

In elderly or poor-performance patients, hypofractionated schedules
(e.g. 40 Gy in 15 fractions) are commonly used as an alternative to 60
Gy in 30 fractions {[}2,3,15{]}. TCP models estimate the probability of
sterilising clonogenic tumour cells as a function of dose and
radiobiological parameters {[}4,5{]}. They are widely used in treatment
planning research but require (i) a well-defined tumour control
endpoint, (ii) inter-patient dose variation, and (iii) sufficient sample
size {[}7{]}.

The CFB-GBM dataset (Centre François Baclesse, n=264) provides pre- and
post-treatment MRI, RTDOSE, GTV contours, clinical covariates,
and---since June 2026---RANO response labels {[}8--10{]}. We asked:
\textbf{can a standard TCP pipeline be validated on this open cohort,
and if not, what does fail and where might signal remain?}

\begin{center}\rule{0.5\linewidth}{0.5pt}\end{center}

\subsection{2. Materials and Methods}\label{materials-and-methods}

\subsubsection{2.1 Dataset and cohort}\label{dataset-and-cohort}

Clinical and imaging metadata were downloaded from TCIA (Version 3, DOI:
\href{https://doi.org/10.7937/v9pn-2f72}{10.7937/v9pn-2f72}) {[}8{]}.
The public CFB-GBM release provides \textbf{pre-processed NIfTI}
(skull-stripped, co-registered) rather than raw DICOM-RT; RTPLAN and
RTSTRUCT DICOM files are therefore not available in the TCIA package. We
used \textbf{RTDOSE} and \textbf{GTV} segmentation masks at t0, which is
the structure set provided for dosimetric analysis. \textbf{CTV and PTV
contours are not included} in this cohort; all DVH metrics refer to GTV
only.

Inclusion required t0 RTDOSE and GTV NIfTI, known fractionation dose and
fraction number. One patient was excluded after DVH QC (Dmean=0 Gy),
yielding \textbf{n=190} for modeling (120×60 Gy/30 fr; 61×40.05 Gy/15
fr) {[}2,15{]}.

\subsubsection{2.2 DVH feature extraction}\label{dvh-feature-extraction}

Cumulative DVH and scalar metrics (D95, Dmean, gEUD, homogeneity index,
GTV volume) were computed from registered RTDOSE and t0 GTV masks. EQD2
was calculated with α/β=10 Gy {[}6,11{]}. gEUD was computed with
volume-effect parameter \emph{a} as in Niemierko {[}6{]}.

\subsubsection{2.3 TCP models}\label{tcp-models}

We implemented four binary-outcome TCP models following standard
radiobiological formulations {[}4,5,7{]}:

\begin{itemize}
\tightlist
\item
  \textbf{Poisson TCP:}
  \(TCP(D) = \exp\!\left[-\ln 2 \cdot \exp\left(\gamma_{50}(1 - D/D_{50})\right)\right]\)
\item
  \textbf{Logistic TCP:} sigmoid in dose with parameters \(D_{50}\),
  \(k\)
\item
  \textbf{Probit TCP:} probit link with \(D_{50}\), \(\sigma\)
\item
  \textbf{gEUD TCP:} dose metric
  \(D = \left(\sum v_i D_i^a\right)^{1/a}\) with \(a \in \{-10,1,10\}\)
  {[}6{]}
\end{itemize}

Parameters were estimated by maximum likelihood. Bootstrap 95\%
confidence intervals (1000 resamples) were computed for Poisson
\(D_{50}\) and \(\gamma_{50}\) {[}18{]}. Models were compared by AIC,
BIC, ROC AUC, and Hosmer--Lemeshow calibration {[}7{]}.

\subsubsection{2.4 Outcomes}\label{outcomes}

\textbf{Primary (exploratory):} OS ≥ cohort median (51 weeks) as a
binary proxy for pipeline testing.

\textbf{Secondary (v3):} RANO non-progressive disease at t1 (non-PD:
SD/MR/PR/CR vs PD) {[}12,13{]}, available for 137/190 modeling patients
{[}8{]}.

\subsubsection{2.5 Statistical analysis}\label{statistical-analysis}

Kaplan--Meier and log-rank tests compared fractionation schemes
{[}14{]}. Cox proportional hazards models included EQD2, Dmean, age,
sex, WHO PS, and RANO non-PD {[}14{]}. Within-arm Spearman correlations
and Poisson TCP AUC were computed separately for 60 Gy and 40 Gy arms.
\textbf{Pooled logistic regression} modeled RANO non-PD
\textasciitilde{} GTV volume + age + WHO PS + scheme indicator, with an
optional volume×scheme interaction term. \textbf{Leave-one-out
cross-validation (LOOCV)} provided out-of-sample AUC for 40 Gy and
pooled models {[}18{]}. \textbf{PyRadiomics features} (v3 TSV, t1gd
sequence, GTV t0) {[}8,9{]} were merged by patient ID; the top five
univariate features by RANO AUC were compared against DVH volume
(in-sample), with \textbf{nested 5-fold stratified CV} (feature
selection on training folds only) to reduce optimism {[}17,18{]}. Full
equations: \texttt{reports/manuscript\_equations\_fragment.tex}
(included in \texttt{manuscript.tex}). Export:
\texttt{bash\ scripts/export\_manuscript.sh} → \texttt{manuscript.docx}
/ \texttt{manuscript.pdf}.

\subsubsection{2.6 Reproducibility}\label{reproducibility}

All code is available at {[}repository URL{]}. Figures and tables
regenerate via \texttt{make\ report}.

\begin{center}\rule{0.5\linewidth}{0.5pt}\end{center}

\subsection{3. Results}\label{results}

\subsubsection{3.1 Cohort and clinical
prognosis}\label{cohort-and-clinical-prognosis}

Median age 70 years; median OS 51 weeks. Sixty Gy patients had longer OS
than 40 Gy patients (median 60 vs 28 weeks; log-rank p≈3×10⁻⁶),
consistent with prior GBM fractionation studies {[}2,3,15{]}. Cox model:
scheme HR≈0.54 (p≈0.0007), WHO PS HR≈1.42 (p≈0.001) {[}14{]}.

\subsubsection{3.2 Pooled TCP (OS proxy)}\label{pooled-tcp-os-proxy}

Poisson/Logistic/Probit TCP on EQD2 achieved in-sample ROC AUC≈0.68
(5-fold CV≈0.68±0.09), significantly better than null (LR p≈3×10⁻⁶).
Bootstrap Poisson: \(D_{50}\)≈53 Gy {[}50, 57{]}, \(\gamma_{50}\)≈3.3
{[}2.1, 4.7{]} {[}18{]}. \textbf{Interpretation:} association is driven
largely by fractionation scheme confounding (r(age, EQD2)=−0.57)
{[}7,15{]}.

\subsubsection{3.3 RANO endpoint (pooled)}\label{rano-endpoint-pooled}

On the same 137 patients, pooled EQD2→RANO non-PD yielded AUC≈0.43 (LR
p=1.0): higher EQD2 correlated with \emph{more} early PD at t1 (60 Gy PD
rate 27\% vs 40 Gy 15\%) {[}12,15{]}.

\subsubsection{3.4 Within-arm analysis}\label{within-arm-analysis}

Within 60 Gy, GTV Dmean SD=0.28 Gy --- insufficient for dose--response
{[}7{]}. \textbf{40 Gy arm (n=34 with RANO):} GTV volume predicted RANO
non-PD (Poisson TCP AUC=0.83, LR p=0.037; Spearman p=0.016), consistent
with tumour burden prognostic literature {[}16{]}. Multivariable
logistic (volume + age + WHO PS): \textbf{AUC=0.90}, bootstrap 95\% CI
{[}0.81, 1.00{]} {[}18{]}.

\subsubsection{3.5 Volume validation}\label{volume-validation}

DVH GTV volume at t0 agreed with RANO \texttt{size\_t0\_cm3} (Spearman
ρ=1.00, n=141) {[}8{]}. t1 NIfTI mask validation was not performed
locally; analyses used author-reported RANO volumes from the v3 TSV
{[}8,9{]}.

\subsubsection{3.6 RANO and overall
survival}\label{rano-and-overall-survival}

Cox model (n=137): RANO non-PD HR≈0.48 (p≈0.0009); EQD2 remained
significant after adjustment (HR≈0.97, p≈0.009) {[}12,14{]}.

\subsubsection{3.7 Pooled RANO models (volume + clinical +
scheme)}\label{pooled-rano-models-volume-clinical-scheme}

On n=137 with RANO labels, EQD2 alone achieved AUC≈0.57. \textbf{GTV
volume + age + WHO PS + scheme} improved discrimination to
\textbf{AUC≈0.72}; adding a volume×scheme interaction yielded AUC≈0.72.
LOOCV AUC for the pooled clinical model was \textbf{≈0.64} {[}18{]},
indicating moderate generalisation within this single-centre cohort.

\subsubsection{3.8 LOOCV in the hypofractionated
arm}\label{loocv-in-the-hypofractionated-arm}

In-sample multivariable AUC in the 40 Gy arm was 0.90, but \textbf{LOOCV
AUC was 0.74} (volume only) and 0.70 (volume + age + PS) {[}18{]},
confirming signal beyond chance while reducing optimism bias relative to
resubstitution {[}7{]}.

\subsubsection{3.9 PyRadiomics vs DVH
volume}\label{pyradiomics-vs-dvh-volume}

Using the author-provided PyRadiomics TSV (t1gd, GTV t0) {[}8,9,17{]},
the top five radiomics features achieved \textbf{in-sample AUC≈0.78} for
RANO non-PD vs \textbf{0.71} for DVH volume alone (n=137).
\textbf{Nested 5-fold CV} (top-5 feature selection on train folds only)
yielded \textbf{AUC≈0.74} for radiomics vs \textbf{0.70} for volume-only
--- optimism Δ≈0.04--0.07. Combined volume + top radiomics feature:
in-sample AUC≈0.77. For OS median-split on the same patients, radiomics
+ clinical covariates reached AUC≈0.78 vs volume-only 0.59 {[}9,10{]}
(see \texttt{figures/08\_pyradiomics\_nested\_cv\_auc.png}).

\subsubsection{3.10 Literature comparison --- TCP parameters (Part
VI)}\label{literature-comparison-tcp-parameters-part-vi}

Poisson TCP on EQD2 (OS median-split proxy) yielded \textbf{D50≈53.2 Gy}
{[}49.5--56.7{]} and \textbf{γ50≈3.32} {[}2.06--4.69{]} (bootstrap,
n=190) {[}18{]}. Direct comparison with published GBM TCP studies is
limited because (i) CFB-GBM lacks formal local control, (ii) dose spread
within protocol is \textless1 Gy, and (iii) our endpoint is OS proxy
rather than LC {[}7,16{]}. Literature review (Table in
\texttt{reports/metrics/literature\_tcp\_d50\_comparison.csv}) shows
reported D50 ranges of \textasciitilde40--80 Gy across sites and
endpoints {[}4,5{]}; our estimate falls within this range but is
\textbf{not interpretable as GBM tumour-control D50} without endpoint
and cohort redesign.

\begin{center}\rule{0.5\linewidth}{0.5pt}\end{center}

\subsection{4. Discussion}\label{discussion}

\subsubsection{4.1 Why pooled dose--TCP fails
here}\label{why-pooled-dosetcp-fails-here}

CFB-GBM reflects \textbf{routine practice}, not a dose-escalation trial:
within-protocol GTV Dmean variation is \textless1 Gy (60 Gy arm SD=0.28
Gy) {[}7,8{]}. TCP fitting further requires a tumour-control endpoint;
OS and early RANO capture different biology and are confounded by
age-linked fractionation selection {[}8,15{]}. Pooled EQD2→RANO AUC≈0.43
is not a failure of RANO per se {[}12,13{]} but of \textbf{using dose as
the sole predictor across mixed fractionation schemes} where higher EQD2
correlates with more early PD at t1 despite better OS {[}15{]}.

Our negative pooled dose--TCP result is scientifically informative: it
demonstrates that \textbf{endpoint and cohort design matter more than
model sophistication} {[}4,7{]}. This aligns with Ohri et
al.\textquotesingle s emphasis on adequate dose heterogeneity and
appropriate endpoints for TCP validation {[}7{]}.

\subsubsection{4.2 Tumour burden predicts RANO when confounding is
addressed}\label{tumour-burden-predicts-rano-when-confounding-is-addressed}

When fractionation scheme is included as a covariate, \textbf{GTV volume
consistently predicts RANO non-PD} (pooled AUC≈0.72; LOOCV≈0.64). This
reframes the analysis from classical TCP (dose→control) to
\textbf{tumour burden→imaging response}, which is clinically
interpretable: larger baseline GTV may reflect more aggressive biology
or reduced likelihood of early imaging stabilisation under
palliative-intent hypofractionation {[}3,15,16{]}.

The hypofractionated subgroup (n=34) shows the strongest in-sample
effect (AUC≈0.90), but \textbf{LOOCV AUC≈0.74} is the more defensible
estimate {[}18{]}. The 60 Gy arm shows a weaker but directionally
consistent association (Spearman p≈0.019, Poisson AUC≈0.66, n=96).

\subsubsection{4.3 PyRadiomics comparison and relation to CFB-GBM
authors}\label{pyradiomics-comparison-and-relation-to-cfb-gbm-authors}

Moreau et al. emphasise imaging AI and radiomics for treatment efficacy
prediction on this cohort {[}9,10{]}. Our head-to-head comparison on the
same open resource {[}8{]} shows that \textbf{author-provided
PyRadiomics features (t1gd GTV t0) modestly outperform DVH scalar
volume} for RANO (in-sample AUC 0.78 vs 0.71; nested CV 0.74 vs 0.70),
while DVH volume remains competitive and fully reproducible from
RTDOSE+GTV without MRI preprocessing.

Feature standardisation follows IBSI recommendations {[}17{]}; we did
not re-extract radiomics locally, ensuring parity with the published
TCIA feature table {[}8,9{]}.

\subsubsection{4.4 Clinical prognostic
context}\label{clinical-prognostic-context}

Independent of TCP framing, CFB-GBM confirms known prognostic factors:
60 vs 40 Gy OS separation (p≈10⁻⁷) {[}2,15{]}, WHO PS gradient {[}2{]},
and RANO non-PD as an OS predictor (Cox HR≈0.48) {[}12,14{]}. These
descriptive findings support data quality but are not novel; they anchor
the dosimetry/RANO feasibility narrative.

\subsubsection{4.5 Implications for open-data TCP
research}\label{implications-for-open-data-tcp-research}

We propose a practical checklist for future open-cohort TCP studies
{[}7{]}: (1) report within-arm GTV Dmean IQR before fitting TCP; (2)
prefer imaging response or LC over OS {[}12,13{]}; (3) adjust for
fractionation and age {[}15{]}; (4) report LOOCV or nested CV alongside
in-sample AUC {[}18{]}; (5) compare against author-provided radiomics
baselines when available {[}8,9,17{]}.

\subsubsection{4.6 Limitations}\label{limitations}

\begin{itemize}
\tightlist
\item
  No formal local control labels; RANO at t1 is an imaging surrogate
  {[}12,13{]}
\item
  Single institution, retrospective, deceased-patients-only inclusion
  {[}8{]}
\item
  Small hypofractionated RANO subset (n=34); pooled LOOCV AUC≈0.64
  limits clinical translation {[}7,18{]}
\item
  t1 GTV NIfTI follow-up masks not validated locally;
  \texttt{size\_t1\_cm3} taken from v3 TSV {[}8{]}
\item
  PyRadiomics nested CV (top-5 selection on train) reduces but does not
  eliminate optimism (Δ≈0.04--0.07 vs in-sample) {[}17,18{]}
\item
  No external validation cohort
\end{itemize}

\begin{center}\rule{0.5\linewidth}{0.5pt}\end{center}

\subsection{5. Conclusion}\label{conclusion}

We deliver an open, reproducible TCP modeling pipeline on CFB-GBM
{[}8{]} and show that \textbf{classical pooled TCP dose--response
validation is not feasible} on this cohort {[}4,7{]}. When fractionation
confounding is addressed {[}15{]}, \textbf{GTV volume and PyRadiomics
features predict early RANO} with moderate discrimination (pooled LOOCV
AUC≈0.64; PyRadiomics nested CV AUC≈0.74 vs volume 0.70) {[}9,17{]}. We
recommend future TCP studies prioritise single-protocol cohorts with ≥1
Gy GTV Dmean IQR, explicit tumour control endpoints {[}12,13{]}, LOOCV
reporting {[}18{]}, and comparison against author radiomics baselines
{[}8,9{]}.

\textbf{LaTeX / Word export:}
\texttt{bash\ scripts/export\_manuscript.sh} →
\texttt{reports/manuscript.tex}, \texttt{manuscript.docx},
\texttt{manuscript.pdf}

\begin{center}\rule{0.5\linewidth}{0.5pt}\end{center}

\subsection{References}\label{references}

\begin{enumerate}
\def\labelenumi{\arabic{enumi}.}
\item
  Stupp R, Mason WP, van den Bent MJ, et al. Effects of radiotherapy
  with concomitant and adjuvant temozolomide versus radiotherapy alone
  on survival in glioblastoma. \emph{Lancet Oncol}. 2009;10(5):459-466.
  https://doi.org/10.1016/S1470-2045(09)70025-7
\item
  Minniti G, De Sanctis V, Valeriani MC. Radiotherapy for glioblastoma:
  current standards and recent advances. \emph{Expert Rev Anticancer
  Ther}. 2021;21(5):529-542.
  https://doi.org/10.1080/14737140.2021.1919470
\item
  Malmström A, Sørensen P, Grunnet K, et al. Temozolomide versus
  standard 6-week radiotherapy versus hypofractionated radiotherapy in
  patients older than 60 years with glioblastoma: the Nordic randomised,
  phase 3 trial. \emph{Lancet Oncol}. 2012;13(9):916-926.
  https://doi.org/10.1016/S1470-2045(12)70265-6
\item
  Maitre A, Vogin G, Kraus R, et al. Construction of radiobiological
  models as TCP and NTCP. \emph{Cancer Radiother}. 2020;24(3):247-257.
  https://doi.org/10.1016/j.canrad.2019.08.002
\item
  Gardner LL, Parry A, McMahon SJ, et al. Modelling radiobiology.
  \emph{Phys Med Biol}. 2024;69(18):18TR01.
  https://doi.org/10.1088/1361-6560/ad70f0
\item
  Niemierko A. Reporting and analyzing dose distributions: a concept of
  equivalent uniform dose. \emph{Med Phys}. 1997;24(1):103-110.
  https://doi.org/10.1118/1.598061
\item
  Ohri N, Tomé WA, Mendez Romero A, et al. Increasing the power of
  tumour control and normal tissue complication probability modelling.
  \emph{Transl Cancer Res}. 2017;6(S1):S123-S127.
  https://doi.org/10.21037/tcr.2017.02.01
\item
  Moreau NN, Leclercq AG, Desmonts A, et al. Pre and post treatment MRI
  and radiotherapy plans of patients with glioblastoma: the CFB-GBM
  cohort (Version 3). \emph{The Cancer Imaging Archive}. 2025.
  https://doi.org/10.7937/v9pn-2f72
\item
  Moreau NN, Valable S, Jaudet C, et al. Early characterization and
  prediction of glioblastoma treatment efficacy using radiomics and AI.
  \emph{Front Oncol}. 2025;15:1497195.
  https://doi.org/10.3389/fonc.2025.1497195
\item
  Moreau NN, Desmonts A, Jaudet C, et al. AI-Driven Prediction of
  Treatment Efficacy in Glioblastoma Using Medical Imaging. In: Rekik I,
  et al., eds. \emph{PRIME 2025}, LNCS vol. 16164. Springer; 2025:24-33.
  https://doi.org/10.1007/978-3-032-07904-6\_3
\item
  Fowler JF. 21 years of biologically effective dose. \emph{Br J
  Radiol}. 2010;83(991):554-568. https://doi.org/10.1259/bjr/86359321
\item
  Wen PY, Macdonald DR, Reardon DA, et al. Updated response assessment
  criteria for high-grade gliomas (RANO). \emph{J Clin Oncol}.
  2010;28(11):1963-1972. https://doi.org/10.1200/JCO.2009.25.5874
\item
  van Dijk WT, van Erp WG, Wesseling P, et al. RANO 2.0 update for
  response assessment in neuro-oncology. \emph{Lancet Oncol}.
  2021;22(12):e503-e508. https://doi.org/10.1016/S1470-2045(21)00489-7
\item
  Cox DR. Regression models and life-tables. \emph{J R Stat Soc Series
  B}. 1972;34(2):187-220.
  https://doi.org/10.1111/j.2517-6161.1972.tb00999.x
\item
  Perry JR, Laperriere N, O\textquotesingle Callaghan CJ, et al.
  Short-course radiation plus temozolomide in elderly patients with
  glioblastoma. \emph{N Engl J Med}. 2017;376(11):1027-1037.
  https://doi.org/10.1056/NEJMoa1611977
\item
  Embring A, Nilsson P, Zellman P, et al. Dosimetric parameters and
  survival in glioblastoma treated with chemoradiation. \emph{Radiother
  Oncol}. 2020;142:47-53. https://doi.org/10.1016/j.radonc.2019.11.014
\item
  Zwanenburg A, Vallières M, Abdalah MA, et al. The Image Biomarker
  Standardisation Initiative. \emph{Radiother Oncol}. 2020;142:169-176.
  https://doi.org/10.1016/j.radonc.2019.11.009
\item
  Efron B, Tibshirani RJ. \emph{An Introduction to the Bootstrap}.
  Chapman \& Hall; 1994. ISBN 978-041204231-5.
\end{enumerate}

\begin{center}\rule{0.5\linewidth}{0.5pt}\end{center}

\subsection{Figure list (from
repository)}\label{figure-list-from-repository}

\begin{longtable}[]{@{}
  >{\raggedright\arraybackslash}p{(\linewidth - 2\tabcolsep) * \real{0.4706}}
  >{\raggedright\arraybackslash}p{(\linewidth - 2\tabcolsep) * \real{0.5294}}@{}}
\toprule\noalign{}
\begin{minipage}[b]{\linewidth}\raggedright
Figure
\end{minipage} & \begin{minipage}[b]{\linewidth}\raggedright
Content
\end{minipage} \\
\midrule\noalign{}
\endhead
\bottomrule\noalign{}
\endlastfoot
Fig 1 & Cohort / fractionation KM
(\texttt{04\_clinical\_prognosis.png}) \\
Fig 2 & Pooled TCP curves OS proxy
(\texttt{03\_tcp\_curves\_os\_proxy.png}) \\
Fig 3 & OS vs RANO AUC comparison
(\texttt{05\_rano\_vs\_os\_tcp\_auc.png}) \\
Fig 4 & Within-arm DVH→RANO (\texttt{06\_within\_arm\_rano\_tcp.png}) \\
Fig 5 & Multivariable ROC 40 Gy
(\texttt{07\_rano\_logistic\_roc\_40gy.png}) \\
Fig 6 & Volume validation scatter
(\texttt{07\_rano\_volume\_validation\_40gy.png}) \\
Fig 7 & Pooled RANO ROC (\texttt{08\_pooled\_rano\_roc.png}) \\
Fig 8 & PyRadiomics vs volume AUC
(\texttt{08\_pyradiomics\_vs\_volume\_auc.png}) \\
Fig 9 & PyRadiomics nested CV
(\texttt{08\_pyradiomics\_nested\_cv\_auc.png}) \\
\end{longtable}

\begin{center}\rule{0.5\linewidth}{0.5pt}\end{center}

\subsection{Presentation outline (15
slides)}\label{presentation-outline-15-slides}

\begin{enumerate}
\def\labelenumi{\arabic{enumi}.}
\tightlist
\item
  Title \& motivation (TCP + open data) {[}4,8{]}
\item
  CFB-GBM overview (TCIA v3, RANO) {[}8--10{]}
\item
  Cohort flowchart (264→190)
\item
  Methods pipeline diagram
\item
  TCP model equations (1 slide) {[}4--6,11{]}
\item
  Clinical results: OS by scheme + WHO PS {[}2,15{]}
\item
  Pooled TCP OS proxy (AUC 0.68) + caveat {[}7{]}
\item
  RANO pooled failure (AUC 0.43) --- confounding {[}12,15{]}
\item
  Pooled volume+scheme models (AUC 0.72) + LOOCV 0.64
\item
  Within-arm 40 Gy: volume→RANO + LOOCV 0.74 {[}16{]}
\item
  PyRadiomics vs DVH volume (AUC 0.78 vs 0.71) {[}9,17{]}
\item
  Cox: RANO predicts OS {[}14{]}
\item
  Limitations (honest)
\item
  Comparison to Moreau radiomics + checklist {[}9,10{]}
\item
  Conclusions + DOI {[}8{]}
\end{enumerate}

\begin{center}\rule{0.5\linewidth}{0.5pt}\end{center}

\emph{Numbers verified against \texttt{reports/metrics/} on 2026-06-28.
Reference DOIs verified in \texttt{reports/literature\_table.csv}.}

\end{document}
